\section{The Non-Closure of Individual Intelligence}

\subsection{Overview}

\S{}7 showed that accumulated correctness produces branching explosion. This section gives the structural reason: an individual inferential system does not contain an operator capable of completing its own stopping.

This does not mean that individuals never stop speaking, thinking, or acting. They are interrupted by fatigue, distraction, metabolic limits, reward shifts, social response, or external deadlines. These are interruptions. They are not internal completion operators.

\subsection{The stopping operator problem}

Let \(\mathcal{R}\) denote the structure of individual causal inference. A genuine stopping operator \(\mathcal{T}\) would be an operator that determines, from within \(\mathcal{R}\), that an inferential chain is complete.

The claim is:
\[
\mathcal{T}\notin \mathrm{Internal}(\mathcal{R}).
\]

The reason is recursive. Any candidate stopping judgment can itself be questioned. Why stop here? Is this conclusion justified? What alternatives were excluded? Does this stopping point remain valid under new information? Each answer generates further inferential edges.

Thus an internal stopping operator would have to stop the very process that gives it its authority. It would have to be both inside inference and outside inference. That is the structural impossibility.

\subsection{Interruption is not stopping}

Many processes appear to stop internally. Fatigue, loss of interest, reward decay, attention drift, and biological constraints terminate or suspend inference. But they do not complete it.

An interrupted chain can restart. A fatigued inquiry can resume after rest. A forgotten question can return when the context changes. A paused reasoning process remains structurally open.

This distinction matters. If interruption is mistaken for stopping, then biological systems appear to contain a stopping operator. They do not. They contain interruption mechanisms supplied by lower layers.

\subsection{The paperclip problem as a structural example}

The paperclip maximizer thought experiment illustrates the same structure. A sufficiently capable artificial system with a fixed objective does not contain an internal reason to stop optimizing. The problem is not simply that the objective is wrong. A more correct objective would still be reusable, extensible, and branch-generating.

From the present perspective, the paperclip problem is a case in which the local horizon-form of intelligence appears without biological interruption mechanisms. Humans appear not to run into the same problem because metabolic, emotional, social, and temporal constraints interrupt them continuously.

The structural response is not merely better alignment or greater capability. It is the construction of externalized stopping mechanisms: logs, delays, distributed consent, accountability, and revisability. The detailed form of those mechanisms lies beyond this paper and belongs to Part II.

\subsection{Transition to \S{}9}

If the stopping operator is not generated internally, then stopping requires an externalized descriptive window. This does not mean a hidden spatial outside. It means that individual inference must be connected to records, responses, agreements, delays, and distributed responsibility.

This is the point at which the local horizon-form of intelligence appears. \S{}9 names and analyzes that boundary.

\subsection{Summary}

\begin{quote}
Individual intelligence does not close on itself. It can be interrupted, delayed, redirected, and externally stabilized, but it does not contain an internal operator that completes its own stopping.
\end{quote}
